\section{Conclusions}
\label{sec:conclusion}

We introduced a reconstruction pipeline whose supervision is a compact, analytic, frozen
2D Gaussian field per view instead of the source photographs.  After a one-time fit per
view the optimization never touches the images again, and supervision amounts to point
queries against kilobyte-scale fields.  The dataset-scale RGB working set thus disappears
from the reconstruction process.  The mechanism is implemented end to end with a verified
point-rendering parity anchor, an unbiased sampled objective and a hard structural
no-image boundary.  \todo{The memory and parity result awaits its controlled experiment
under the preregistered protocol of \cref{sec:memory}.}

We also showed that the same captures admit a closed-form tomographic initialization,
since rendering projects 3D Gaussians to 2D Gaussians and the projection can be inverted.
Beam Fusion combines back-projected beams by covariance intersection, requires three
contributing views due to an observability argument and yields complete contributor
lineage.  The carrier refinement repairs the identifiable failures of the naive inversion
with stages justified by paired ablations.  \todo{Whether the tomographic initialization
beats the standard initializations under matched budgets is an open, preregistered
question, and a negative answer leaves the memory result untouched.}

Building on the standard 3DGS optimizer and renderer makes our approach directly
applicable to existing pipelines.  The natural next step is temporal, from 2D Gaussian
video over space-time beam fusion to a 4D carrier representation for standard 4D Gaussian
splatting.
